\documentclass[11pt]{article}
\usepackage{amsmath,amssymb,amsthm}
\usepackage{hyperref}
\usepackage[margin=1in]{geometry}

\newcommand{\N}{\mathbb{N}}
\newcommand{\Z}{\mathbb{Z}}
\newcommand{\Q}{\mathbb{Q}}
\newcommand{\GK}{\operatorname{GK}}
\newcommand{\CPP}{\operatorname{CPP}}
\newcommand{\IP}{\operatorname{IP}}
\newcommand{\ALCD}{\operatorname{ALCD}}
\newcommand{\qpoch}[2]{(#1;#2)_\infty}
\newcommand{\qpochn}[3]{(#1;#2)_{#3}}
\newcommand{\qbin}[2]{\genfrac{[}{]}{0pt}{}{#1}{#2}}

\theoremstyle{plain}
\newtheorem{theorem}{Theorem}
\newtheorem{conjecture}[theorem]{Conjecture}
\newtheorem{proposition}[theorem]{Proposition}
\newtheorem{lemma}[theorem]{Lemma}
\newtheorem{corollary}[theorem]{Corollary}
\theoremstyle{definition}
\newtheorem{definition}[theorem]{Definition}
\newtheorem{example}[theorem]{Example}
\newtheorem{remark}[theorem]{Remark}

\title{The Underlying Set: Cylindric Plane Partitions, Arbitrarily Labelled
Cylindric Diagrams, and the Ground Set for the Involution}
\author{Notes by Robin Langer}
\date{\today}

\begin{document}
\maketitle

\begin{abstract}
We give a detailed account of the bijection between cylindric plane
partitions and pairs consisting of an integer partition and an
arbitrarily labelled cylindric diagram (ALCD).  This bijection,
constructed via cylindric growth diagrams, is the combinatorial
incarnation of the $E$--$H$ factorisation of the generating function
$\GK_c(z,q)$.  We then use it to give a precise description of the
underlying set $\mathcal{G}_{n,c}$ on which the sign-reversing
involution for the positivity conjecture must act.
\end{abstract}

\tableofcontents

%% ============================================================
\section{The bijection}
\label{sec:bijection}
%% ============================================================

\subsection{Statement}

The main result of the PhD thesis (Chapter~3) is a weight-preserving
bijection:
\begin{equation}\label{eq:bijection}
  \psi_\pi \colon (\IP,\; \ALCD(\pi)) \xrightarrow{\;\sim\;} \CPP(\pi),
\end{equation}
where:
\begin{itemize}
  \item $\IP$ is the set of all integer partitions (the \emph{base
    partition} $\gamma$).
  \item $\ALCD(\pi)$ is the set of arbitrarily labelled cylindric
    diagrams of profile~$\pi$.
  \item $\CPP(\pi)$ is the set of cylindric plane partitions of
    profile~$\pi$.
\end{itemize}

The bijection decomposes each CPP into two independent pieces: an
ordinary partition $\gamma$ (the ``universal'' piece, independent of
the profile) and an ALCD $\mathfrak{d}$ (the ``purely cylindric'' piece,
which carries all the information about the profile and the interlacing
conditions).

\subsection{Algebraic counterpart: the $E$--$H$ factorisation}

At the level of generating functions, the bijection corresponds to the
factorisation
\begin{equation}\label{eq:EH}
  \GK_c(z,q) = \frac{1}{\qpoch{zq}{q}} \cdot G_c(z,q),
\end{equation}
where $H(z) := 1/\qpoch{zq}{q} = \sum_{n \geq 0} q^n z^n/\qpochn{q}{q}{n}$
is the generating function for ordinary partitions (with largest part
$\leq n$, or equivalently with at most~$n$ parts, at $z$-degree~$n$;
see Section~\ref{sec:why-z-max}), and
$G_c(z,q) := \qpoch{zq}{q} \cdot \GK_c(z,q)$ is the corrected
generating function encoding the purely cylindric content.

\begin{center}
\begin{tabular}{l|l}
\textbf{Algebraic} & \textbf{Bijective} \\ \hline
$\GK_c(z,q) = H(z) \cdot G_c(z,q)$ &
  $\CPP(c) \xleftrightarrow{\;\psi\;} \IP \times \ALCD(c)$ \\[3pt]
$H(z) = 1/\qpoch{zq}{q}$ &
  The base partition $\gamma$ \\[3pt]
$G_c(z,q)$ &
  The ALCD $\mathfrak{d}$
\end{tabular}
\end{center}


%% ============================================================
\section{Why $z$ tracks the maximum entry}
\label{sec:why-z-max}
%% ============================================================

The generating function $\GK_c(z,q) = \sum_{\mathfrak{c}} z^{\max(\mathfrak{c})}\,
q^{|\mathfrak{c}|}$ uses $z$ to track the maximum entry
$\max(\mathfrak{c}) = \max_k \mu^k_1$ (the largest first part among all
slices).  Why is this the natural statistic?

\subsection{The conjugation symmetry for partitions}

For a single partition $\lambda$, the \emph{largest part} $\lambda_1$ and
the \emph{number of parts} $\ell(\lambda)$ are exchanged by conjugation:
$\lambda_1 = \ell(\lambda')$ and $\ell(\lambda) = \lambda'_1$.  Since
conjugation is a size-preserving involution $|\lambda| = |\lambda'|$,
the two statistics are \emph{equidistributed}:
\begin{equation}\label{eq:conjugation-symmetry}
  \sum_\lambda z^{\lambda_1}\, q^{|\lambda|}
  = \sum_\lambda z^{\ell(\lambda)}\, q^{|\lambda|}
  = \frac{1}{\qpoch{zq}{q}}.
\end{equation}

\begin{proof}
For the second equality, expand the product:
$\prod_{k \geq 1} 1/(1-zq^k)$.
The factor $1/(1-zq^k) = \sum_{a \geq 0} z^a q^{ak}$ contributes
multiplicity~$a$ for parts of size~$k$, with $z^a$ tracking the count.
So $z$ tracks the total number of parts $\ell(\lambda) = \sum_k a_k$.

For the first equality, substitute $\lambda \mapsto \lambda'$ in the
left-hand sum: $\sum_\lambda z^{\lambda_1} q^{|\lambda|}
= \sum_{\lambda'} z^{\ell(\lambda')} q^{|\lambda'|}
= \sum_\mu z^{\ell(\mu)} q^{|\mu|}$.
\end{proof}

This means the universal factor $H(z) = 1/\qpoch{zq}{q}$ can be read
in two equivalent ways:
\begin{center}
\begin{tabular}{l|l}
\textbf{Reading} & \textbf{$z^n$ counts} \\ \hline
$z$ = largest part ($\lambda_1 = n$)
  & Partitions with largest part exactly~$n$ \\[2pt]
$z$ = number of parts ($\ell(\lambda) = n$)
  & Partitions with exactly~$n$ parts
\end{tabular}
\end{center}

\subsection{The two geometric interpretations}

In the growth diagram, these two readings correspond to two
complementary geometric pictures:

\paragraph{Height interpretation ($z = $ largest part).}
The max entry $\max(\mathfrak{c})$ measures the \emph{height} of
the tallest column in any slice of the CPP.  This is the statistic
that appears in the definition of $\GK_c(z,q)$ and in the conjecture:
extracting $[z^n]\,\GK_c(z,q)$ selects CPPs whose tallest column has
height exactly~$n$.

\paragraph{Depth interpretation ($z = $ number of parts).}
The growth diagram has a natural \emph{depth} parameter: the number of
peeling layers required to reduce the CPP to its base partition.  Each
layer removes entries at one height level (from $m$ down to~$1$).  A CPP
with max entry~$m$ requires exactly~$m$ peeling layers.  At each layer,
the face labels of the ALCD record the data extracted.

The conjugation symmetry~\eqref{eq:conjugation-symmetry} guarantees that
these two pictures are consistent: a partition with largest part~$m$
(height~$m$) produces, after conjugation, a partition with $m$~parts
(depth~$m$).

\subsection{The $E$--$H$ factorisation and the power of $z$}

For the full CPP (not just the base partition), the $E$--$H$ factorisation
\[
  \GK_c(z,q) = \frac{1}{\qpoch{zq}{q}} \cdot G_c(z,q)
\]
decomposes $z^{\max(\mathfrak{c})}$ as a \emph{convolution}: extracting
$[z^n]$ gives
\begin{equation}\label{eq:convolution}
  f_{n,c}(q) = [z^n]\,\GK_c(z,q)
  = \sum_{j=0}^{n} \frac{q^j}{\qpochn{q}{q}{j}} \cdot g_{n-j,c}(q),
\end{equation}
where $q^j/\qpochn{q}{q}{j}$ counts partitions $\gamma$ with
$\ell(\gamma) = j$ (equivalently, $\gamma_1 = j$), and $g_{n-j,c}(q)$
is the corrected coefficient from $G_c(z,q)$.

At the level of the bijection, this convolution reflects the fact that
the max entry of a CPP receives contributions from both the base
partition~$\gamma$ and the ALCD~$\mathfrak{d}$:
\[
  \max(\mathfrak{c}) = \underbrace{\gamma_1}_{\text{base partition height}}
  \;+\; \underbrace{(\max(\mathfrak{c}) - \gamma_1)}_{\text{ALCD contribution}},
\]
where the ALCD contribution is the number of additional height levels
added by the forward local rules beyond the base partition.

\begin{remark}
The corrected expansion $H(z) = \sum_{n \geq 0} q^n z^n/\qpochn{q}{q}{n}$
(not $\sum z^n/\qpochn{q}{q}{n}$, which would be $1/\qpoch{z}{q}$)
reflects the fact that a partition with exactly~$n$ parts has size at
least~$n$ (each part $\geq 1$), contributing the factor $q^n$.
\end{remark}


%% ============================================================
\section{Cylindric diagrams}
\label{sec:alcd}
%% ============================================================

\subsection{The cylindric poset}

\begin{definition}[Cylindric poset]
\label{def:cyl-poset}
For integers $n, m \geq 0$ with $T := n + m > 0$, the
\emph{cylindric poset} $\mathfrak{G}(n,m)$ is the quotient of $\Z^2$
by the equivalence relation $(x, y) \sim (x + n,\, y - m)$.  Its
elements are equivalence classes of lattice points, inheriting the
partial order from $\Z^2$.
\end{definition}

A binary string $\pi = (\pi_0, \pi_1, \ldots, \pi_{T-1}) \in \{0,1\}^T$
with $n$ ones and $m$ zeros encodes the profile: $\pi_k = 1$ means
``add a horizontal strip'' (going up) and $\pi_k = 0$ means ``remove
a horizontal strip'' (going down).

For rank~3 with profile $c = (c_0, c_1, c_2)$ and $d = c_0 + c_1 + c_2$,
the binary string $\pi$ has length $T = d$ and is determined by the
partition $\mu(c) = (c_1 + c_2,\, c_2,\, 0)$ associated to the profile.


\subsection{Arbitrarily labelled cylindric diagrams}

\begin{definition}[ALCD]
\label{def:alcd}
An \emph{arbitrarily labelled cylindric diagram} (ALCD) of profile $\pi$
is an assignment of non-negative integers (``face labels'') to the faces
of the cylindric poset $\mathfrak{G}(n,m)$, such that all but finitely
many face labels are zero.

More precisely, an ALCD $\mathfrak{d}$ is a function
\[
  \mathfrak{d}\colon \{\text{faces of } \mathfrak{G}(n,m)\} \to \Z_{\geq 0}
\]
with finite support.  The \emph{weight} of $\mathfrak{d}$ is
$|\mathfrak{d}| := \sum_f \mathfrak{d}(f)$, the sum of all face labels.
\end{definition}

\begin{remark}[Why ``arbitrarily labelled'']
The face labels are \textbf{arbitrary} non-negative integers --- there
are no inequalities between adjacent face labels.  This is in contrast
to the vertex labels in a growth diagram, which must satisfy interlacing
conditions.  The word ``arbitrary'' refers to this freedom.
\end{remark}

\begin{remark}[Diagonal structure]
\label{rem:diagonals}
The faces of $\mathfrak{G}(n,m)$ are naturally grouped into
\emph{diagonals}.  For rank~3 with profile $c = (c_0, c_1, c_2)$, the
faces at each ``depth level'' $k$ form a diagonal with $T = d$ faces
(one for each position around the cylinder).  A face at depth~$k$ and
position~$i$ is denoted $\mathfrak{d}(k, i)$, where $k \geq 1$ and
$i \in \{0, 1, \ldots, T-1\}$ (cyclic).

The diagonal weight at depth~$k$ is
$|\mathrm{diag}(k)|_{\mathfrak{d}} := \sum_{i=0}^{T-1} \mathfrak{d}(k,i)$.
\end{remark}


%% ============================================================
\section{The growth diagram bijection}
\label{sec:growth-diagram}
%% ============================================================

\subsection{Growth diagrams}

\begin{definition}[Cylindric growth diagram]
\label{def:growth-diagram}
A \emph{cylindric growth diagram} of profile $\pi \in \{0,1\}^T$ is a
labelling of the cylindric poset $\mathfrak{G}(n,m)$ consisting of:
\begin{itemize}
  \item \textbf{Vertex labels:} an integer partition $\lambda(v)$ at
    each vertex~$v$.
  \item \textbf{Face labels:} a non-negative integer $\mathfrak{d}(f)$
    at each face~$f$.
\end{itemize}
These must satisfy:
\begin{enumerate}
  \item If $v \lhd w$ (covering relation) then $\lambda(v)/\lambda(w)$
    or $\lambda(w)/\lambda(v)$ is a horizontal strip (depending on
    whether the edge is ``up'' or ``down'').
  \item All but finitely many vertices carry the same partition~$\gamma$
    (the \emph{base partition}).
  \item At each face with vertices $(\mu; \alpha, \beta; \lambda)$
    and face label~$m$, the \emph{local rule} is satisfied:
    $\lambda = \mathfrak{U}_{\alpha,\beta}(m, \mu)$.
\end{enumerate}
\end{definition}

The growth diagram interpolates between the two sides of the bijection:
\begin{itemize}
  \item The \textbf{upper boundary} (the outermost ring of vertex
    labels) gives a CPP $\mathfrak{c} = (\mu^0, \mu^1, \ldots,
    \mu^{T-1})$.
  \item The \textbf{face labels} give an ALCD $\mathfrak{d}$.
  \item The \textbf{interior} converges to the base partition $\gamma$
    at sufficient depth.
\end{itemize}


\subsection{The Burge local rule}

The local rule $\mathfrak{U}_{\alpha,\beta}$ is the engine that drives
the bijection.  It has type signature
\begin{equation}\label{eq:local-rule-type}
  \mathfrak{U}_{\alpha,\beta}\colon
  \Z_{\geq 0} \times (D(\alpha) \cap D(\beta))
  \to U(\alpha) \cap U(\beta),
\end{equation}
where $U(\mu)$ (resp.\ $D(\lambda)$) is the set of partitions obtained
by adding (resp.\ removing) a horizontal strip to $\mu$
(resp.\ from $\lambda$).

\begin{definition}[Burge forward rule]
\label{def:burge-forward}
Given partitions $\alpha$, $\beta$, $\nu$ (with $\nu \leq \alpha$ and
$\nu \leq \beta$ in the dominance order induced by horizontal strips)
and a non-negative integer~$m$, define $\lambda =
\mathfrak{U}_{\alpha,\beta}(m, \nu)$ by:
\begin{enumerate}
  \item Work in conjugate coordinates: $\alpha' = \mathrm{conj}(\alpha)$,
    $\beta' = \mathrm{conj}(\beta)$, $\nu' = \mathrm{conj}(\nu)$.
  \item Compute the column sets:
    \begin{align*}
      A &:= \{i \geq 1 : \alpha'_i > \nu'_i\}
        && \text{(columns where $\alpha$ exceeds $\nu$)}, \\
      B &:= \{i \geq 1 : \beta'_i > \nu'_i\}
        && \text{(columns where $\beta$ exceeds $\nu$)}.
    \end{align*}
  \item \textbf{Cascade:} for each $i \in A \cap B$ (in increasing
    order), let $\varepsilon(i)$ be the smallest integer $> i$ not
    in $A \cup B$ and not previously assigned.
    Set $C := \{\varepsilon(i) : i \in A \cap B\}$.
  \item \textbf{Free columns:} let $D$ be the first $m$ elements of
    $\Z_{>0} \setminus (A \cup B \cup C)$ (in increasing order).
  \item \textbf{Build $\lambda$:} set $\lambda'_i := \nu'_i + [i \in A \cup B \cup C \cup D]$
    (add one box to each target column).
\end{enumerate}
\end{definition}

\begin{definition}[Burge inverse rule]
\label{def:burge-inverse}
Given partitions $\alpha$, $\beta$, $\lambda$ (with $\lambda \geq \alpha$
and $\lambda \geq \beta$), the inverse rule recovers $(m, \nu) =
\mathfrak{U}_{\alpha,\beta}^{-1}(\lambda)$ by:
\begin{enumerate}
  \item Work in conjugate coordinates.
  \item Compute the column sets:
    \begin{align*}
      \bar{A} &:= \{i \geq 1 : \lambda'_i > \alpha'_i\}, \\
      \bar{B} &:= \{i \geq 1 : \lambda'_i > \beta'_i\}.
    \end{align*}
  \item \textbf{Reverse cascade:} for each $i \in \bar{A} \cap \bar{B}$
    (in \emph{decreasing} order), let $\delta(i)$ be the largest
    integer $< i$ not in $\bar{A} \cup \bar{B}$ and not previously
    assigned.
    Set $\bar{C} := \{\delta(i) : \delta(i) > 0\}$ and
    $m := \#\{\delta(i) \leq 0\}$.
  \item \textbf{Build $\nu$:} set $\nu'_i := \lambda'_i -
    [i \in \bar{A} \cup \bar{B} \cup \bar{C}]$.
\end{enumerate}
\end{definition}

\begin{proposition}[Weight condition]
\label{prop:weight}
The local rule satisfies the weight identity:
\begin{equation}\label{eq:weight-cond}
  |\lambda| - |\nu| = |A \cup B \cup C \cup D|
  = |A| + |B| - |A \cap B| + |C| + m,
\end{equation}
where $|C| = |A \cap B|$, so $|\lambda| - |\nu| = |A| + |B| + m$.
In the growth diagram context (where $\alpha/\nu$ and $\beta/\nu$ are
horizontal strips), this gives
\begin{equation}\label{eq:weight-growth}
  |\lambda| = |\alpha| + |\beta| - |\nu| + m.
\end{equation}
\end{proposition}


\subsection{Constructing the bijection}

Given a pair $(\gamma, \mathfrak{d}) \in \IP \times \ALCD(\pi)$, the
CPP $\psi_\pi(\gamma, \mathfrak{d})$ is constructed by:
\begin{enumerate}
  \item \textbf{Initialise:} label all vertices at sufficient depth
    with the base partition $\gamma$.
  \item \textbf{Build upward:} at each face, apply the forward local
    rule $\mathfrak{U}_{\alpha,\beta}(m, \nu) = \lambda$, using the
    face label $m = \mathfrak{d}(f)$ and the vertex labels already
    determined.
  \item \textbf{Read off:} the upper boundary gives the CPP
    $(\mu^0, \mu^1, \ldots, \mu^{T-1})$.
\end{enumerate}

Conversely, given a CPP $\mathfrak{c}$, the pair $(\gamma, \mathfrak{d})$
is recovered by:
\begin{enumerate}
  \item \textbf{Initialise:} label the upper boundary with the CPP
    partitions.
  \item \textbf{Peel downward:} at each face, apply the inverse rule
    $\mathfrak{U}_{\alpha,\beta}^{-1}(\lambda) = (m, \nu)$, extracting
    the face label $m$ and the lower vertex label $\nu$.
  \item \textbf{Converge:} after finitely many layers, all vertex labels
    stabilise at the base partition $\gamma$.  The extracted face labels
    constitute the ALCD $\mathfrak{d}$.
\end{enumerate}


\subsection{Strong weight preservation}

The bijection is not merely weight-preserving but
\textbf{strongly} weight-preserving.

\begin{proposition}[PhD thesis, Chapter~3]
\label{prop:strong-weight}
If $\psi_\pi(\gamma, \mathfrak{d}) = (\mu^0, \mu^1, \ldots, \mu^{T-1})$,
then for all $0 \leq k \leq T-1$:
\begin{equation}\label{eq:strong-weight}
  |\mu^k| = |\gamma| + |\mathrm{diag}(k)|_{\mathfrak{d}},
\end{equation}
where $|\mathrm{diag}(k)|_{\mathfrak{d}}$ is the total weight of the
$k$-th diagonal of the ALCD (Remark~\ref{rem:diagonals}).

In particular, the total weight of the CPP is
\[
  |\mathfrak{c}| = \sum_{k=0}^{T-1} |\mu^k|
  = T \cdot |\gamma| + |\mathfrak{d}|,
\]
where $|\mathfrak{d}| = \sum_k |\mathrm{diag}(k)|_{\mathfrak{d}}$
is the total weight of the ALCD.
\end{proposition}

This strong form is essential: it means the bijection preserves not
just the total weight but the individual slice sizes.  This is the
property that makes the bijection suitable for analysing the
\emph{multivariate} generating function (which tracks each $|\mu^k|$
separately).


\subsection{The max statistic under the bijection}

For the bivariate generating function $\GK_c(z,q) = \sum z^{\max(\mathfrak{c})} q^{|\mathfrak{c}|}$, the max statistic of a CPP
$\mathfrak{c} = (\mu^0, \ldots, \mu^{T-1})$ is
\[
  \max(\mathfrak{c}) = \max_{0 \leq k \leq T-1} \mu^k_1
\]
(the largest first part among all slices).

Under the bijection $\psi_\pi(\gamma, \mathfrak{d}) = \mathfrak{c}$,
the max statistic depends on \emph{both} $\gamma$ and $\mathfrak{d}$:
\[
  \max(\mathfrak{c}) = \gamma_1 + \max_{k} \bigl(\text{max column index
  used at depth~1, position~$k$}\bigr).
\]
More precisely, the first part $\mu^k_1$ is determined by the base
partition $\gamma$ plus the effect of all face labels in the $k$-th
column of the growth diagram.  The exact formula involves the cascade
mechanism of the local rule.

\begin{remark}
The non-trivial dependence of the max statistic on both $\gamma$ and
$\mathfrak{d}$ is why the $E$--$H$ factorisation at the level of
$\GK_c(z,q)$ is more subtle than at the level of $\GK_c(1,q)$: the
$z$-grading (by max) cannot be expressed as a simple product of
$z$-gradings on $\gamma$ and $\mathfrak{d}$ separately.
\end{remark}


%% ============================================================
\section{The polynomials $h_{m,c}(q)$ as ALCD generating functions}
\label{sec:hmc}
%% ============================================================

\subsection{Factoring out the base partition}

Recall the key generating functions:
\begin{align}
  f_{m,c}(q) &:= [z^m]\,\GK_c(z,q)
    = \sum_{\substack{\mathfrak{c} \in \CPP(c) \\ \max(\mathfrak{c}) = m}}
      q^{|\mathfrak{c}|},
    \label{eq:fmc} \\
  h_{m,c}(q) &:= \qpochn{q}{q}{m} \cdot f_{m,c}(q).
    \label{eq:hmc}
\end{align}

The factor $1/\qpochn{q}{q}{m}$ generates partitions with largest part
$\leq m$ (equivalently, with at most $m$ parts after conjugation).
So $f_{m,c}(q)$ counts CPPs with $\max = m$, and division by
$1/\qpochn{q}{q}{m}$ amounts to ``quotienting out'' the contribution of
the base partition $\gamma$ (restricted to have $\gamma_1 \leq m$).

Under the bijection, a CPP $\mathfrak{c}$ with $\max(\mathfrak{c}) = m$
decomposes as $(\gamma, \mathfrak{d})$ where $\gamma$ contributes
$q^{|\gamma|}$ to the weight and $\mathfrak{d}$ contributes
$q^{|\mathfrak{d}|}$.  The multiplication by $\qpochn{q}{q}{m}$
``cancels'' the sum over $\gamma$, leaving:

\begin{equation}\label{eq:hmc-alcd}
  h_{m,c}(q) = \sum_{\substack{\mathfrak{d} \in \ALCD(c) \\
  \max(\psi(\gamma_0, \mathfrak{d})) = m}}
  q^{|\mathfrak{d}|},
\end{equation}
where $\gamma_0$ is the unique base partition such that the resulting
CPP has $\max = m$.

\begin{remark}[Why this is subtle]
The constraint $\max(\psi(\gamma, \mathfrak{d})) = m$ couples $\gamma$
and $\mathfrak{d}$: the max of the CPP depends on both the base
partition and the ALCD.  The factorisation
$f_{m,c}(q) = h_{m,c}(q)/\qpochn{q}{q}{m}$ is exact at the level of
generating functions, but the combinatorial interpretation of the
individual objects counted by $h_{m,c}(q)$ requires specifying how
the max constraint is distributed between $\gamma$ and $\mathfrak{d}$.

Nevertheless, computational evidence confirms that $h_{m,c}(q) \in
\Z_{\geq 0}[q]$ for all tested profiles.  This is a strictly stronger
statement than positivity of $Q_{n,c}(q)$.
\end{remark}


%% ============================================================
\section{The ground set for the involution}
\label{sec:ground-set}
%% ============================================================

We now bring together all the pieces to give a precise description of
the underlying set on which the sign-reversing involution must act.

\subsection{The polynomial $Q_{n,c}(q)$}

\begin{definition}[$Q_{n,c}(q)$]
\label{def:Q}
For a nonneg integer $n$ and profile $c = (c_0, c_1, c_2)$ with
$d = c_0 + c_1 + c_2$ and $\ell = \gcd(d, 3)$, define
\begin{equation}\label{eq:Q-def}
  Q_{n,c}(q) := \qpochn{q^\ell}{q^\ell}{n}
  \cdot [z^n]\!\Bigl(\qpoch{zq}{q} \cdot \GK_c(z,q)\Bigr)
  = \qpochn{q^\ell}{q^\ell}{n} \cdot g_{n,c}(q).
\end{equation}
When $d \not\equiv 0 \pmod{3}$, we have $\ell = 1$ and
$Q_{n,c}(q) = \qpochn{q}{q}{n} \cdot g_{n,c}(q)$.
\end{definition}

\subsection{The $q$-binomial decomposition}

Recall the \emph{$q$-binomial coefficient}
\[
  \qbin{n}{j} := \frac{\qpochn{q}{q}{n}}{\qpochn{q}{q}{j}\,\qpochn{q}{q}{n-j}},
\]
which is a polynomial in $q$ equal to the generating function for
partitions fitting inside a $j \times (n{-}j)$ rectangle:
$\qbin{n}{j} = \sum_{\lambda \subseteq (n-j)^j} q^{|\lambda|}$.

The polynomial $Q_{n,c}(q)$ has the alternating-sum representation:
\begin{equation}\label{eq:Q-alt}
  Q_{n,c}(q)
  = \sum_{j=0}^{n} (-1)^j\, q^{j(j+1)/2}\,
    \qbin{n}{j} \cdot h_{n-j,c}(q).
\end{equation}

Using the identity $q^{j(j+1)/2} \qbin{n}{j} = \sum_{\substack{S
\subseteq \{1,\ldots,n\} \\ |S| = j}} q^{\sigma(S)}$
(where $\sigma(S) = \sum_{s \in S} s$), this becomes:
\begin{equation}\label{eq:Q-ground}
  Q_{n,c}(q)
  = \sum_{\substack{S \subseteq \{1,\ldots,n\} \\
    \tau \in \mathcal{H}_{n-|S|,c}}}
    (-1)^{|S|}\, q^{\sigma(S) + \operatorname{wt}(\tau)},
\end{equation}
where $\mathcal{H}_{m,c}$ is the set of objects counted by $h_{m,c}(q)$,
equipped with the weight function $\operatorname{wt}$.

\begin{remark}[From the analytic definition to the ground set]
\label{rem:analytic-to-ground}
The analytic definition of $Q_{n,c}(q)$ (see the companion note on the
positivity conjecture) is
\begin{equation}\label{eq:Q-analytic}
  Q_{n,c}(q) = \qpochn{q}{q}{n}
  \cdot [z^n]\!\bigl(\qpoch{zq}{q} \cdot \GK_c(z,q)\bigr).
\end{equation}
(Here we state the case $d \not\equiv 0 \pmod{3}$, i.e.\ $\ell =
\gcd(d,3) = 1$; the general case replaces $(q)_n$ by
$(q^\ell;q^\ell)_n$.)
Let us trace through the three steps that convert this into the
ground-set expression~\eqref{eq:Q-ground}.

\medskip
\noindent\textbf{Step~1: Expand $\qpoch{zq}{q}$ via Euler's identity.}
The Euler identity gives
$\qpoch{zq}{q} = \sum_{j \geq 0} (-1)^j q^{j(j+1)/2} z^j /
\qpochn{q}{q}{j}$.
Writing $f_{m,c}(q) = [z^m]\,\GK_c(z,q)$ for the coefficients of the
original generating function, extracting $[z^n]$ yields
\[
  [z^n]\!\bigl(\qpoch{zq}{q} \cdot \GK_c(z,q)\bigr)
  = \sum_{j=0}^{n} \frac{(-1)^j\, q^{j(j+1)/2}}{\qpochn{q}{q}{j}}
    \cdot f_{n-j,c}(q)
  = g_{n,c}(q).
\]

\noindent\textbf{Step~2: Multiply by $\qpochn{q}{q}{n}$.}
Using $\qpochn{q}{q}{n}/\qpochn{q}{q}{j} =
\qpochn{q}{q}{n-j} \cdot \qbin{n}{j}$ and the definition
$h_{m,c}(q) = \qpochn{q}{q}{m} \cdot f_{m,c}(q)$, we obtain the
$q$-binomial decomposition~\eqref{eq:Q-alt}:
\[
  Q_{n,c}(q) = \sum_{j=0}^{n} (-1)^j\, q^{j(j+1)/2}\,
  \qbin{n}{j} \cdot h_{n-j,c}(q).
\]
This is the step where the denominators are cleared: each
$1/\qpochn{q}{q}{j}$ is absorbed into a $q$-binomial coefficient,
ensuring $Q_{n,c}(q)$ is a polynomial (not merely a formal power series).

\medskip
\noindent\textbf{Step~3: Biject $q$-binomials to subsets.}
The following identity converts each $q$-binomial into a sum over
subsets, producing the ground-set expression~\eqref{eq:Q-ground}.
\end{remark}

\begin{lemma}\label{lem:qbin-subsets}
For integers $0 \leq j \leq n$,
\begin{equation}\label{eq:qbin-subsets}
  q^{j(j+1)/2}\,\qbin{n}{j}
  = \sum_{\substack{S \subseteq \{1,\ldots,n\} \\ |S| = j}} q^{\sigma(S)},
\end{equation}
where $\sigma(S) = \sum_{s \in S} s$.
\end{lemma}

\begin{proof}
There is a bijection between $j$-element subsets
$S = \{s_1 < \cdots < s_j\} \subseteq \{1,\ldots,n\}$ and partitions
$\lambda$ fitting inside a $j \times (n{-}j)$ rectangle, given by
\[
  \lambda_i = s_{j+1-i} - (j + 1 - i), \qquad i = 1,\ldots, j.
\]
Since $s_i \geq i$ (there are at least $i$ elements of $S$ that are
$\leq s_i$), each $\lambda_i \geq 0$.  The strict ordering
$s_1 < \cdots < s_j$ gives
$\lambda_1 \geq \lambda_2 \geq \cdots \geq \lambda_j \geq 0$,
so $\lambda$ is a partition with at most $j$ parts.  Moreover
$s_j \leq n$ gives $\lambda_1 = s_j - j \leq n - j$, so
$\lambda$ fits inside the box $(n{-}j)^j$.

The weight satisfies
\[
  |\lambda| = \sum_{i=1}^{j}(s_i - i) = \sigma(S) - \tfrac{j(j+1)}{2},
\]
so $q^{\sigma(S)} = q^{|\lambda| + j(j+1)/2}$.  Summing over all
$j$-element subsets and using
$\qbin{n}{j} = \sum_{\lambda \subseteq (n-j)^j} q^{|\lambda|}$:
\[
  \sum_{|S|=j} q^{\sigma(S)}
  = q^{j(j+1)/2} \sum_{\lambda \subseteq (n-j)^j} q^{|\lambda|}
  = q^{j(j+1)/2}\,\qbin{n}{j}. \qedhere
\]
\end{proof}


\subsection{The ground set: three equivalent descriptions}

We give three descriptions of the ground set, at increasing levels
of explicitness.

\medskip
\noindent\textbf{Level 1: Abstract pairs.}
The ground set is
\begin{equation}\label{eq:G-abstract}
  \mathcal{G}_{n,c} =
  \bigsqcup_{j=0}^{n}
  \bigl\{(S, \tau) : S \subseteq \{1,\ldots,n\},\;
    |S| = j,\; \tau \in \mathcal{H}_{n-j,c}\bigr\},
\end{equation}
with sign $\varepsilon(S,\tau) = (-1)^{|S|}$ and weight
$w(S,\tau) = \sigma(S) + \operatorname{wt}(\tau)$.

Here $\tau$ is an \emph{abstract} element of the set
$\mathcal{H}_{m,c}$ counted by $h_{m,c}(q)$.  The advantage of this
description is its simplicity; the disadvantage is that $\tau$ is not a
concrete combinatorial object.

\medskip
\noindent\textbf{Level 2: Pairs with ALCDs.}
Using equation~\eqref{eq:hmc-alcd}, we can identify
$\tau \in \mathcal{H}_{m,c}$ with an ALCD $\mathfrak{d}$ at the
appropriate max level:
\begin{equation}\label{eq:G-alcd}
  \mathcal{G}_{n,c} =
  \bigsqcup_{j=0}^{n}
  \bigl\{(S, \mathfrak{d}) : S \subseteq \{1,\ldots,n\},\;
    |S| = j,\; \mathfrak{d} \in \ALCD_{n-j}(c)\bigr\},
\end{equation}
where $\ALCD_m(c)$ denotes the set of ALCDs of profile~$c$ that produce
a CPP with $\max = m$ when combined with the appropriate base partition.
The weight is $w(S, \mathfrak{d}) = \sigma(S) + |\mathfrak{d}|$.

Here $\mathfrak{d}$ is a concrete combinatorial object (a finitely-supported
function from faces to $\Z_{\geq 0}$), but the constraint linking it
to ``max level~$m$'' is implicit and involves the base partition.

\medskip
\noindent\textbf{Level 3: Triples with CPPs.}
Expanding $h_{m,c}(q) = \qpochn{q}{q}{m} \cdot f_{m,c}(q)$
and $f_{m,c}(q) = \sum_{\nu:\max(\nu)=m} q^{|\nu|}$, we get the
fully explicit ground set:
\begin{equation}\label{eq:G-triples}
  \widetilde{\mathcal{G}}_{n,c} =
  \bigsqcup_{j=0}^{n}
  \left\{(S, U, \nu) :
    \begin{array}{l}
      S \subseteq \{1,\ldots,n\},\; |S| = j, \\[2pt]
      U \subseteq \{1,\ldots,n-j\}, \\[2pt]
      \nu \in \CPP(c),\; \max(\nu) = n - j
    \end{array}
  \right\},
\end{equation}
with sign $\varepsilon(S, U, \nu) = (-1)^{|S| + |U|}$ and weight
$w(S, U, \nu) = \sigma(S) + \sigma(U) + |\nu|$.

Here $U$ comes from the factor $\qpochn{q}{q}{m}$ and $\nu$ is a
genuine CPP.  The expression for $Q_{n,c}(q)$ is:
\begin{equation}
  Q_{n,c}(q) = \sum_{(S, U, \nu) \in \widetilde{\mathcal{G}}_{n,c}}
  (-1)^{|S|+|U|}\, q^{\sigma(S)+\sigma(U)+|\nu|}.
\end{equation}


\subsection{The relationship between the three levels}

The three descriptions are related by:
\[
  \text{Level 3 (triples)}
  \;\xrightarrow{\text{cancel $U$ vs.\ $\nu$}}\;
  \text{Level 2 (ALCD pairs)}
  \;\xrightarrow{\text{forget structure}}\;
  \text{Level 1 (abstract pairs)}.
\]

The passage from Level~3 to Level~2 uses the bijection
$\psi_\pi\colon (\gamma, \mathfrak{d}) \leftrightarrow \nu$ (with
$\gamma$ absorbed into $U$ via the identity
$h_{m,c} = (q)_m \cdot f_{m,c}$).  Since $h_{m,c}(q)$ is
(computationally verified to be) nonneg, the sign from $(-1)^{|U|}$
cancels internally, and we can work at Level~1 or Level~2 where only
$(-1)^{|S|}$ contributes to the sign.


\subsection{Summary: what the involution must do}

A sign-reversing, weight-preserving involution
$\iota\colon \mathcal{G}_{n,c} \to \mathcal{G}_{n,c}$ must satisfy:
\begin{enumerate}
  \item $\iota^2 = \operatorname{id}$.
  \item Non-fixed points pair with opposite sign: if
    $\iota(S, \tau) = (S', \tau') \neq (S, \tau)$,
    then $|S'| \equiv |S| + 1 \pmod{2}$.
  \item Weight is preserved:
    $\sigma(S) + \operatorname{wt}(\tau)
    = \sigma(S') + \operatorname{wt}(\tau')$.
  \item Every fixed point $(S, \tau) = \iota(S, \tau)$ has
    $|S| \equiv 0 \pmod{2}$.
\end{enumerate}

The fixed-point generating function gives $Q_{n,c}(q)$:
\[
  Q_{n,c}(q) = \sum_{\substack{(S,\tau) \in \mathcal{G}_{n,c} \\
  \iota(S,\tau) = (S,\tau)}} q^{\sigma(S) + \operatorname{wt}(\tau)}.
\]

At Level~2 (ALCD pairs), the peel/glue mechanism of the local rules
provides the natural transformation engine: the \textbf{inverse local
rule} extracts face labels from the ALCD (reducing its weight) and
deposits them into the subset~$S$, while the \textbf{forward local rule}
does the reverse.  The challenge is to define a \emph{canonical} choice
of which face label to peel or glue, determined entirely by the pair
$(S, \mathfrak{d})$, such that the resulting map is involutive.


%% ============================================================
\section{Worked example}
\label{sec:example}
%% ============================================================

\subsection{Setup: $d = 4$, $c = (2,2,0)$, $n = 1$}

For this case, $Q_{1,(2,2,0)}(q) = q^2 + q^3 + q^4 + q^5$.

\paragraph{The $h$-polynomials.}
\begin{itemize}
  \item $h_{0,c}(q) = 1$ (the empty ALCD).
  \item $h_{1,c}(q) = q + q^2 + q^3 + q^4 + q^5$
    (five ALCDs at max level~1, with weights $1, 2, 3, 4, 5$).
\end{itemize}

\paragraph{The alternating sum.}
\[
  Q_{1,c}(q) = h_{1,c}(q) - q \cdot h_{0,c}(q)
  = (q + q^2 + q^3 + q^4 + q^5) - q = q^2 + q^3 + q^4 + q^5.
\]

\paragraph{The ground set $\mathcal{G}_{1,c}$ at Level~1.}
\begin{itemize}
  \item $j = 0$ (sign $+$): pairs $(\emptyset, \tau)$ with
    $\tau \in \mathcal{H}_{1,c}$.
    Five objects at weights $1, 2, 3, 4, 5$.
  \item $j = 1$ (sign $-$): pairs $(\{1\}, \tau_0)$ with
    $\tau_0 \in \mathcal{H}_{0,c}$.
    One object at weight $\sigma(\{1\}) + 0 = 1$.
\end{itemize}

Total: six objects, five positive and one negative.

\paragraph{The involution.}
The unique negative object $(\{1\}, \tau_0)$ at weight~$1$ pairs with
the unique positive object $(\emptyset, \tau_1^{\min})$ at weight~$1$
(where $\tau_1^{\min}$ is the minimal-weight ALCD at max level~1).

Bijectively, this pairing says: the ALCD $\tau_1^{\min}$ of weight~$1$
can be ``peeled'' (via an inverse local rule that extracts a face label
of value~$1$) to produce the empty ALCD $\tau_0$ of weight~$0$.

The four surviving fixed points at weights $2, 3, 4, 5$ give
$Q_{1,c}(q) = q^2 + q^3 + q^4 + q^5$. \qed


%% ============================================================
\section{Outlook: what the bijection provides}
\label{sec:outlook}
%% ============================================================

The growth diagram bijection
$\psi_\pi\colon (\IP, \ALCD(\pi)) \to \CPP(\pi)$ provides:

\begin{enumerate}
  \item \textbf{A concrete model for $\mathcal{H}_{m,c}$.}
    The objects $\tau$ counted by $h_{m,c}(q)$ are (equivalence classes
    of) ALCDs at the appropriate max level.

  \item \textbf{A transformation mechanism.}
    The forward and inverse local rules provide the peel/glue operations
    needed for the involution: the inverse rule extracts face labels
    (peeling), the forward rule deposits them (gluing).

  \item \textbf{Weight tracking.}
    The strong weight preservation (Proposition~\ref{prop:strong-weight})
    ensures that peel/glue operations preserve the total $q$-weight,
    which is the essential requirement for the involution.

  \item \textbf{Multivariate compatibility.}
    The strong weight preservation tracks individual slice sizes, so the
    same framework should extend to the multivariate conjecture.
\end{enumerate}

The missing ingredient is the \textbf{canonical choice}: which face
label to peel or glue, and how to decide between peeling and gluing,
given a pair $(S, \mathfrak{d})$.  This choice must be globally
consistent (involutive) and must produce fixed points only at
even $|S|$.

\begin{thebibliography}{99}
\bibitem{Borodin}
  A.~Borodin,
  \emph{Periodic Schur process and cylindric partitions},
  Duke Math.~J.\ \textbf{140} (2007), 391--468.

\bibitem{Fomin}
  S.~V.~Fomin,
  \emph{Generalized Robinson--Schensted--Knuth correspondence},
  J.\ Soviet Math.\ \textbf{41} (1988), 979--991.

\bibitem{Krattenthaler}
  C.~Krattenthaler,
  \emph{Growth diagrams, and increasing and decreasing chains in
  fillings of Ferrers shapes},
  Adv.\ Appl.\ Math.\ \textbf{37} (2006), 404--431.

\bibitem{Langer-thesis}
  R.~Langer,
  \emph{Cylindric Partitions and Bijective Proofs},
  PhD thesis, University of Melbourne.

\bibitem{Langer-cpp1}
  R.~Langer,
  \emph{Enumeration of Cylindric Plane Partitions -- Part~I},
  2012. arXiv:1204.4583.
\end{thebibliography}

\end{document}
